\section{Separating the boundary from the log canonical cycle}
\label{sec:boundary}

The numerical relation $N_0\sim m(K_X+B)$ initially controls the log
canonical divisor, whereas Theorem~\ref{thm:obstruction} requires a
relation for a power of the canonical bundle. We now show that the
boundary contributes only components of the cycle or divisors entirely
disjoint from it. The argument uses maximal nef dimension and
bend-and-break, and is independent of the subsequent MMP.

\begin{proposition}[Boundary separation]\label{prop:separation}
Let $X$ be a normal $\Q$-factorial projective threefold over an
uncountable algebraically closed field, and let $B$ be an effective
$\Q$-divisor. Suppose that $N$ is a nonzero effective nef Cartier
divisor with $\nu(N)=1$, $n(N)=3$, and
$N\sim_{\Q}a(K_X+B)$ for some rational $a>0$. For every ample Cartier
divisor $H$, one has
\begin{equation}\label{eq:boundary-zero}
BNH=0.
\end{equation}
Consequently $N$ is numerically trivial on each prime component of $B$,
and every such component is either contained in $\Supp N$ or disjoint
from $\Supp N$.
\end{proposition}
\begin{proof}
We may replace $H$ by a very ample multiple. The number $BNH$ is
nonnegative. If it were positive, the identity
$(K_X+B)NH=N^2H/a=0$ would give $K_XNH<0$.

For each sufficiently large integer $l$, choose positive integers
$u_l,v_l$ such that $u_lH$ and $v_l(H+lN)$ are very ample. Their general
complete intersection is an integral curve $C_l$ contained in the
smooth locus of $X$. Indeed, normality gives a singular locus of
dimension at most one; the first general member meets this locus in
finitely many points, and the second avoids them. Bertini's theorem on
the smooth open gives the desired integral curve. Its intersections
satisfy
\begin{equation}\label{eq:bend-ratio}
\frac{NC_l}{-K_XC_l}
=\frac{NH^2}{-K_XH^2-lK_XHN}\longrightarrow0,
\end{equation}
and $-K_XC_l>0$ for large $l$.

The quantitative bend-and-break theorem of Jovinelly, Lehmann, and
Riedl \cite[Theorem~1.1]{JLR} applies to an integral curve in the smooth
locus of a projective variety over an algebraically closed field, with
negative canonical degree, and any nef Cartier test divisor. In
dimension three it supplies, through every closed point of $C_l$, a
rational curve $\Gamma$ satisfying
\[
0\leq N\Gamma\leq4\frac{NC_l}{-K_XC_l}.
\]
Fix $l$ so large that the right-hand side is strictly less than one.
Since $N$ is Cartier, it follows that $N\Gamma=0$.

For completeness, these complete intersections pass through points to
which the very-general-point condition applies. For the two fixed very
ample systems, the incidence of a point of $X$ and a pair of members
through it is a product of projective bundles over $X$, hence is
irreducible and dominates $X$. Restricting to the nonempty open set of
pairs giving the curves just considered leaves a dense open incidence
and a dominant evaluation map. Its image is constructible and contains
a nonempty open subset of $X$. Over the uncountable field, this open
contains a closed point outside the countable union excluded by
$n(N)=3$. The curve $\Gamma$ through that point is a contradiction.
This proves \eqref{eq:boundary-zero}.

Write $B=\sum b_TT$ with $b_T>0$. All numbers $NHT$ are nonnegative,
so \eqref{eq:boundary-zero} gives $NHT=0$ for each $T$.
Lemma~\ref{lem:surface-nef} yields $N|_T\equiv0$. If $T$ is not a
component of $N$, the canonical section of $\mathcal O_X(N)$ restricts
to a nonzero section on the integral surface $T$. Were $T\cap\Supp N$
nonempty, its zero scheme would be a nonzero effective Cartier divisor
on $T$ and would have positive ample degree. This contradicts
$N|_T\equiv0$, so the two supports are disjoint.
\end{proof}

Apply Proposition~\ref{prop:separation} to $N_0$ and set
\begin{equation}\label{eq:Pout}
P_{\mathrm{out}}=
\bigcup_{\substack{T\text{ a prime component of }B\\T\not\subset\Supp N_0}}T.
\end{equation}
This is a closed subset disjoint from $\Supp N_0$. Thus
$\Supp B\subset\Supp N_0\cup P_{\mathrm{out}}$, a union of two
disjoint closed sets. We retain this separation while isolating one
component of $N_0$.
